% Colocar dentro de \chapter{Resultados}. Requer booktabs e \citeonline (já usados no texto).
% Os números das tabelas ainda devem ser preenchidos com a nova execução (comando \pend).

\providecommand{\pend}{\textbf{[a preencher]}}

\chapter{Ablação}
\label{chap:res-ablacao}

Neste capítulo são detalhados os experimentos de ablação conduzidos para determinar as arquiteturas e os hiperparâmetros da via topológica e da via visual antes da composição do modelo de dupla via. As duas vias foram avaliadas de forma isolada sobre as mesmas partições de dados, com o propósito de selecionar as melhores configurações sem utilizar o conjunto de teste final reservado. A organização conceitual da avaliação das vias isoladas está ilustrada na \ref{fig:ablacao-vias}, abrangendo os modelos avaliados na via topológica e na via visual, com suas respectivas cabeças de classificação e extração de métricas.

\begin{figure}[!hbtp]
    \centering
    \resizebox{\columnwidth}{!}{%
    \begin{tikzpicture}[
        main_block/.style={
            rectangle, rounded corners=3pt, inner sep=5pt, align=center,
            font=\small\bfseries, draw=black!70, fill=blue!10, line width=0.8pt
        },
        title_block/.style={
            rectangle, rounded corners=3pt, inner sep=3.5pt, align=center,
            font=\footnotesize\bfseries, draw=black!70, fill=blue!10, line width=0.8pt
        },
        arrow/.style = {thick, ->, >={Stealth[length=2.5mm]}},
        box_delimitador/.style = {
            rectangle, rounded corners=8pt, draw=black!70, fill=black!2, line width=0.8pt
        },
        model/.style = {
            rectangle, rounded corners=3pt, align=center, draw=black!70, line width=0.7pt,
            font=\footnotesize, minimum width=2.5cm, minimum height=0.75cm, inner sep=3pt
        },
        legend_sq/.style = {
            rectangle, draw=black!70, line width=0.5pt, minimum size=0.28cm, inner sep=0pt
        },
        cnn/.style = {fill=blue!15},
        trf/.style = {fill=green!15},
        hyb/.style = {fill=orange!22},
        rec/.style = {fill=blue!15},
        base/.style = {fill=black!10}
    ]

        %%%%%%%%%%%%%%%% (a) VIA TOPOLÓGICA %%%%%%%%%%%%%%%%
        \node (ds_a)  [main_block, minimum width=1.9cm] at (0,0) {\textit{Datasets}};
        \node (ssn_a) [main_block, minimum width=2.2cm] at (3.1,0) {SSN};

        \node (grp_a) [box_delimitador, minimum width=9.0cm, minimum height=3.6cm] at (9.6,0) {};
        \node (grp_a_t) [title_block] at (grp_a.north) {6 modelos testados individualmente};
        \node [model, rec,  minimum width=2.5cm] at (6.9,0.55)  {LSTM};
        \node [model, rec,  minimum width=2.5cm] at (9.6,0.55)  {BiLSTM};
        \node [model, rec,  minimum width=2.5cm] at (12.3,0.55) {GRU};
        \node [model, rec,  minimum width=2.5cm] at (6.9,-0.45) {BiGRU};
        \node [model, rec,  minimum width=2.5cm] at (9.6,-0.45) {SAC-LSTM};
        \node [model, base, minimum width=2.5cm] at (12.3,-0.45) {MLP};
        \node [legend_sq, rec]  at (6.2,-1.35) {};
        \node [anchor=west, font=\scriptsize] at (6.4,-1.35) {Recorrente};
        \node [legend_sq, base] at (8.9,-1.35) {};
        \node [anchor=west, font=\scriptsize] at (9.1,-1.35) {\textit{Baseline} não recorrente};

        \node (head_a) [main_block, minimum width=2.4cm, align=center] at (16.2,0) {Cabeça\\MLP};
        \node (met_a)  [main_block, minimum width=1.9cm] at (19.3,0) {Métricas};

        \draw [arrow] (ds_a.east)  -- (ssn_a.west);
        \draw [arrow] (ssn_a.east) -- (grp_a.west);
        \node [below=2pt of ssn_a, font=\scriptsize] {$15\times18$};
        \draw [arrow] (grp_a.east) -- (head_a.west);
        \draw [arrow] (head_a.east) -- (met_a.west);

        \node [font=\small] at (9.65,-2.35) {(a) Via topológica};

        %%%%%%%%%%%%%%%% (b) VIA VISUAL %%%%%%%%%%%%%%%%
        \begin{scope}[yshift=-6.4cm]
        \node (ds_b)  [main_block, minimum width=1.9cm] at (0,0) {\textit{Datasets}};
        \node (aug_b) [main_block, minimum width=2.2cm, align=center] at (3.1,0) {\footnotesize Aumento de\\[-1pt]\footnotesize dados};

        \node (grp_b) [box_delimitador, minimum width=9.0cm, minimum height=3.6cm] at (9.6,0) {};
        \node (grp_b_t) [title_block] at (grp_b.north) {6 modelos testados individualmente};
        \node [model, cnn] at (6.9,0.55)  {ResNet-50};
        \node [model, cnn] at (6.9,-0.45) {EfficientNet-B0};
        \node [model, trf] at (9.6,0.55)  {ViT-Tiny};
        \node [model, trf] at (9.6,-0.45) {DeiT-Tiny};
        \node [model, hyb] at (12.3,0.55) {MobileViT-S};
        \node [model, hyb] at (12.3,-0.45){CoAtNet-0};
        \node [legend_sq, cnn] at (5.9,-1.35) {};
        \node [anchor=west, font=\scriptsize] at (6.1,-1.35) {Convolucional};
        \node [legend_sq, trf] at (8.4,-1.35) {};
        \node [anchor=west, font=\scriptsize] at (8.6,-1.35) {\textit{Transformer}};
        \node [legend_sq, hyb] at (11.0,-1.35) {};
        \node [anchor=west, font=\scriptsize] at (11.2,-1.35) {Híbrida};

        \node (head_b) [main_block, minimum width=2.4cm, align=center] at (16.2,0) {Cabeça\\MLP};
        \node (met_b)  [main_block, minimum width=1.9cm] at (19.3,0) {Métricas};

        \draw [arrow] (ds_b.east)  -- (aug_b.west);
        \draw [arrow] (aug_b.east) -- (grp_b.west);
        \node [below=2pt of aug_b, font=\scriptsize] {$256\times256$};
        \draw [arrow] (grp_b.east) -- (head_b.west);
        \draw [arrow] (head_b.east) -- (met_b.west);

        \node [font=\small] at (9.65,-2.35) {(b) Via visual};
        \end{scope}

    \end{tikzpicture}%
    }
    \caption{Ablação das vias isoladas. (a) Via topológica: a sequência de características da SSN é processada por seis codificadores (cinco recorrentes e um MLP como \textit{baseline}), cada um testado de forma independente. (b) Via visual: as imagens, após aumento de dados, são processadas por seis modelos (convolucionais, \textit{transformers} e híbridos), também testados individualmente. Cada via possui sua própria cabeça MLP, da qual são extraídas as métricas.}
    \label{fig:ablacao-vias}
\end{figure}


\section{Ablação da Via Topológica}
\label{sec:via-topologica}
\label{sec:ablacao-topologica}

A via topológica recebe a sequência de características calculada a partir de redes complexas sobre 15 raios de conexão, denotada por $X_{\mathrm{top}} \in \mathbb{R}^{15 \times 18}$, em que cada passo contém 18 medidas descritoras. Para verificar a capacidade de aprendizado sobre essa sequência e avaliar a contribuição da dependência de ordem entre as escalas, foram selecionados cinco modelos de redes neurais recorrentes e um modelo não recorrente como \textit{baseline} (\ref{fig:ablacao-vias}a):

\begin{itemize}
    \item a LSTM padrão \cite{hochreiter1997long,gers2000learning}, modelo com portas de esquecimento, entrada e saída, empregado como referência no processamento de sequências multiescala \cite{lin2022mslstm} e diagnósticos médicos \cite{mathiyalagi2026hclstm},
    \item a BiLSTM \cite{schuster1997bidirectional,graves2005framewise}, que processa os raios nos sentidos direto e inverso para integrar informação de micro e macrotextura, com precedentes na detecção de alterações celulares neoplásicas \cite{priyadarshini2026cnnbilstm} e em modelos de dupla via \cite{yousaf2024twostream},
    \item a GRU \cite{cho2014learning}, que unifica as portas e dispensa o estado de célula separado, resultando em menor custo paramétrico para conter o risco de sobreajuste em conjuntos com volume reduzido de imagens,
    \item a BiGRU \cite{schuster1997bidirectional,cho2014learning}, que combina o fluxo de informação bidirecional entre as escalas à redução paramétrica da unidade GRU,
    \item a SAC-LSTM (\textit{Self-Attention Coupled LSTM}) \cite{shalini2026saclstm}, modelo desenvolvido para classificação histopatológica de câncer de mama, que adiciona uma camada de autoatenção aos estados da LSTM para ponderar dinamicamente os raios de conexão mais discriminativos,
    \item o MLP (\textit{Multi-Layer Perceptron}) \cite{huang2022fusing}, rede de alimentação direta que processa as 270 características de forma estática sobre o vetor achatado, avaliada como linha de base de controle para verificar se a modelagem sequencial traz ganho efetivo frente a uma projeção densa sem recorrência.
\end{itemize}

Todos os seis codificadores produzem um vetor de características com a mesma dimensão oculta, conectado à mesma cabeça de classificação, formada por normalização por camadas, camada linear de 128 neurônios com ativação ReLU, camada de \textit{dropout} e camada linear de saída com projeção para o número de classes de cada tarefa.

\label{subsec:protocolo-ablacao}
O protocolo experimental de particionamento, validação cruzada, busca em grade e seleção de hiperparâmetros dos modelos em ablação está esquematizado na \ref{fig:ablacao-treino}. O procedimento foi executado exclusivamente sobre a partição de desenvolvimento de cada conjunto de dados. Inicialmente, as imagens de cada conjunto foram divididas de forma estratificada, com semente 42, na proporção de 80\% para desenvolvimento e 20\% para teste final, permanecendo a partição de teste totalmente reservada e sem qualquer participação no processo de ablação. Sobre os 80\% de desenvolvimento, foi aplicada uma validação cruzada estratificada em cinco partições ($k = 5$). Para cada modelo e para cada combinação de hiperparâmetros, o treinamento foi conduzido em cinco rodadas independentes, nas quais quatro partições foram destinadas ao ajuste dos pesos e uma partição foi utilizada para a validação, alternando a partição de validação a cada rodada. As 18 medidas descritoras foram padronizadas em cada rodada por escore padrão calculado exclusivamente sobre as partições de treino, prevenindo vazamento de dados para a validação. A grade de hiperparâmetros explorada para cada conjunto de dados na via topológica está sintetizada na Tabela~\ref{tab:espaco-busca-top}.

\begin{figure}[!hbtp]
    \centering
    \resizebox{0.75\columnwidth}{!}{%
    \begin{tikzpicture}[
        recntagle/.style={
            rectangle, rounded corners=3pt, inner sep=5pt, align=center,
            font=\small\bfseries, draw=black!70, fill=blue!10, line width=0.8pt
        },
        title_block/.style={
            rectangle, rounded corners=3pt, inner sep=4.5pt, align=center,
            font=\small\bfseries, draw=black!70, fill=blue!10, line width=0.8pt
        },
        arrow/.style = {thick, ->, >={Stealth[length=2.5mm]}},
        box_delimitador/.style = {
            rectangle, rounded corners=8pt, draw=black!70, fill=black!2, line width=0.8pt
        },
        dashed_box/.style = {
            rectangle, rounded corners=4pt, draw=black!70, dashed, line width=0.8pt,
            fill=white, align=center
        },
        sub_block/.style = {
            rectangle, rounded corners=3pt, inner sep=5pt, align=center,
            draw=black!70, fill=blue!5, line width=0.7pt, font=\small
        },
        fold/.style = {
            rectangle, draw=black!70, line width=0.6pt, minimum width=1.3cm,
            minimum height=0.7cm, font=\small, inner sep=2pt
        },
        legend_sq/.style = {
            rectangle, draw=black!70, line width=0.5pt, minimum size=0.3cm, inner sep=0pt
        }
    ]

        % 1. Partição estratificada
        \node (split_box) [box_delimitador, minimum width=8.6cm, minimum height=2.4cm] at (0,0) {};
        \node (split_title) [title_block] at (split_box.north) {1. Partição estratificada (semente 42)};
        \node (dev) [sub_block, minimum width=5.2cm, minimum height=0.95cm] at (-1.5,-0.25) {Desenvolvimento (80\%)};
        \node (test) [dashed_box, minimum width=1.9cm, minimum height=0.95cm] at (2.9,-0.25) {\footnotesize Teste (20\%)\\\scriptsize reservado};

        % 2. Validação cruzada
        \node (cv_box) [box_delimitador, minimum width=8.6cm, minimum height=2.9cm] at (0,-3.9) {};
        \node (cv_title) [title_block] at (cv_box.north) {2. Validação cruzada ($k=5$)};
        \node (f1) [fold, fill=blue!10]   at (-2.9,-3.8) {F1};
        \node (f2) [fold, fill=blue!10]   at (-1.45,-3.8) {F2};
        \node (f3) [fold, fill=blue!10]   at (0,-3.8) {F3};
        \node (f4) [fold, fill=blue!10]   at (1.45,-3.8) {F4};
        \node (f5) [fold, fill=orange!25] at (2.9,-3.8) {F5};
        \node [legend_sq, fill=blue!10]   at (-3.3,-4.85) {};
        \node [anchor=west, font=\scriptsize] at (-3.1,-4.85) {Treino (4 partições)};
        \node [legend_sq, fill=orange!25] at (-0.1,-4.85) {};
        \node [anchor=west, font=\scriptsize] at (0.1,-4.85) {Validação (1 partição)};

        % 3. Busca de hiperparâmetros
        \node (hp_box) [box_delimitador, minimum width=8.6cm, minimum height=2.4cm] at (0,-7.5) {};
        \node (hp_title) [title_block] at (hp_box.north) {3. Busca de hiperparâmetros};
        \node (c1) [sub_block, minimum width=1.7cm] at (-3.1,-7.2) {Comb. 1};
        \node (c2) [sub_block, minimum width=1.7cm] at (-1.05,-7.2) {Comb. 2};
        \node (c3) [sub_block, minimum width=1.7cm] at (1.0,-7.2) {$\cdots$};
        \node (cx) [sub_block, minimum width=1.7cm] at (3.05,-7.2) {Comb. $X$};
        \node [font=\scriptsize, align=center] at (0,-8.2) {cada combinação é treinada nas 5 rodadas\\(4 partições de treino $\times$ 1 de validação, alternada)};

        % 4. Seleção
        \node (sel_box) [box_delimitador, minimum width=8.6cm, minimum height=2.0cm] at (0,-10.9) {};
        \node (sel_title) [title_block] at (sel_box.north) {4. Seleção da configuração};
        \node (avg) [sub_block, text width=3.2cm] at (-2.3,-11.1) {F1 macro médio por combinação};
        \node (best) [sub_block, text width=2.4cm] at (2.6,-11.1) {Melhor configuração};
        \draw [arrow] (avg.east) -- (best.west);

        % Conexões
        \draw [arrow] (split_box.south) -- (cv_title.north);
        \draw [arrow] (cv_box.south) -- (hp_title.north);
        \draw [arrow] (hp_box.south) -- (sel_title.north);

    \end{tikzpicture}%
    }
    \caption{Protocolo de ablação para o treinamento dos modelos: partição estratificada em 80\% para desenvolvimento e 20\% para teste final (reservado), validação cruzada com cinco partições sobre os 80\%, avaliação de $X$ combinações de hiperparâmetros e seleção da combinação de maior desempenho médio.}
    \label{fig:ablacao-treino}
\end{figure}


\begin{table}[htb]
\centering
\caption{Espaço de busca de hiperparâmetros estratificado por conjunto de dados na via topológica.}
\label{tab:espaco-busca-top}
\label{subsec:res-ajuste}
\small
\resizebox{\textwidth}{!}{%
\begin{tabular}{lccccc}
\toprule
\textbf{\textit{Dataset}} & \textbf{Lote ($B$)} & \textbf{Taxa de aprendizado ($\eta$)} & \textbf{Decaimento ($\lambda$)} & \textbf{\textit{Dropout} ($p$)} & \textbf{Dim. oculta ($h$)} \\
\midrule
UCSB & $\{4, 8, 16\}$   & $\{3\cdot 10^{-4}, 10^{-3}, 3\cdot 10^{-3}, 10^{-2}, 3\cdot 10^{-2}\}$ & $10^{-4}$ & $\{0{,}1, 0{,}3, 0{,}5\}$ & $\{64, 128, 256, 512\}$ \\
CR   & $\{16, 32, 64\}$ & $\{3\cdot 10^{-4}, 10^{-3}, 3\cdot 10^{-3}, 10^{-2}, 3\cdot 10^{-2}\}$ & $10^{-4}$ & $\{0{,}1, 0{,}3, 0{,}5\}$ & $\{64, 128, 256, 512\}$ \\
LG   & $\{16, 32, 64\}$ & $\{3\cdot 10^{-4}, 10^{-3}, 3\cdot 10^{-3}, 10^{-2}, 3\cdot 10^{-2}\}$ & $10^{-4}$ & $\{0{,}1, 0{,}3, 0{,}5\}$ & $\{64, 128, 256, 512\}$ \\
NHL  & $\{16, 32, 64\}$ & $\{3\cdot 10^{-4}, 10^{-3}, 3\cdot 10^{-3}, 10^{-2}, 3\cdot 10^{-2}\}$ & $10^{-4}$ & $\{0{,}1, 0{,}3, 0{,}5\}$ & $\{64, 128, 256, 512\}$ \\
LA   & $\{16, 32, 64\}$ & $\{3\cdot 10^{-4}, 10^{-3}, 3\cdot 10^{-3}, 10^{-2}, 3\cdot 10^{-2}\}$ & $10^{-4}$ & $\{0{,}1, 0{,}3, 0{,}5\}$ & $\{64, 128, 256, 512\}$ \\
\bottomrule
\end{tabular}%
}
\end{table}

O processo de treinamento foi conduzido com a função de perda de entropia cruzada, otimizador AdamW e norma de gradiente limitada ao valor 1,0 para conter oscilações numéricas. O coeficiente de decaimento de pesos foi fixado no valor $\lambda = 10^{-4}$. Em decorrência do tamanho das representações e dos volumes de dados, o treinamento foi delimitado ao teto de 100 épocas, com parada antecipada regulada por uma tolerância de 15 épocas consecutivas sem decrescimento da perda na partição de validação, sendo os pesos da melhor época restaurados ao término de cada rodada.

\label{subsec:res-codificadores}
\label{subsec:res-top-selecionadas}
O desempenho de cada combinação foi computado pela média aritmética do F1 macro obtido nas cinco partições de validação, selecionando a combinação responsável pela maior pontuação média como a configuração de referência de cada arquitetura em cada conjunto de dados. Os desempenhos obtidos pelas configurações selecionadas de cada codificador, acompanhados do respectivo desvio padrão, estão registrados na Tabela~\ref{tab:abl-f1}, enquanto as combinações de hiperparâmetros responsáveis por esses resultados estão discriminadas na Tabela~\ref{tab:abl-top-params}.

\begin{table}[htb]
\centering
\caption{F1 macro (média $\pm$ desvio padrão entre as 5 partições de validação) da via topológica isolada, na melhor combinação de hiperparâmetros de cada modelo. Em negrito, o maior valor médio de cada conjunto de dados.}
\label{tab:abl-f1}
\small
\resizebox{\textwidth}{!}{%
\begin{tabular}{lccccc}
\toprule
\textbf{Codificador} & \textbf{UCSB} & \textbf{CR} & \textbf{LA} & \textbf{LG} & \textbf{NHL} \\
\midrule
LSTM & 0{,}893 $\pm$ 0{,}124 & 0{,}992 $\pm$ 0{,}017 & \textbf{1{,}000} $\pm$ 0{,}000 & 0{,}995 $\pm$ 0{,}011 & 0{,}938 $\pm$ 0{,}043 \\
BiLSTM & 0{,}870 $\pm$ 0{,}109 & 0{,}985 $\pm$ 0{,}034 & \textbf{1{,}000} $\pm$ 0{,}000 & 0{,}995 $\pm$ 0{,}011 & 0{,}956 $\pm$ 0{,}019 \\
GRU & 0{,}892 $\pm$ 0{,}107 & \textbf{0{,}993} $\pm$ 0{,}017 & \textbf{1{,}000} $\pm$ 0{,}000 & 0{,}995 $\pm$ 0{,}011 & 0{,}956 $\pm$ 0{,}049 \\
BiGRU & \textbf{0{,}933} $\pm$ 0{,}061 & 0{,}985 $\pm$ 0{,}034 & \textbf{1{,}000} $\pm$ 0{,}000 & \textbf{0{,}995} $\pm$ 0{,}011 & 0{,}966 $\pm$ 0{,}021 \\
SAC-LSTM & 0{,}867 $\pm$ 0{,}090 & 0{,}977 $\pm$ 0{,}034 & \textbf{1{,}000} $\pm$ 0{,}000 & 0{,}995 $\pm$ 0{,}011 & 0{,}921 $\pm$ 0{,}046 \\
MLP & 0{,}871 $\pm$ 0{,}135 & \textbf{0{,}993} $\pm$ 0{,}017 & \textbf{1{,}000} $\pm$ 0{,}000 & 0{,}995 $\pm$ 0{,}011 & \textbf{0{,}973} $\pm$ 0{,}020 \\
\bottomrule
\end{tabular}%
}
\end{table}

As combinações ótimas de hiperparâmetros responsáveis pelos maiores índices de F1 macro em cada conjunto de dados estão discriminadas na Tabela~\ref{tab:abl-top-params}.

\begin{table}[htb]
\centering
\caption{Melhores combinações de hiperparâmetros obtidas na validação cruzada para cada modelo nos cinco conjuntos de dados.}
\label{tab:abl-top-params}
\label{tab:abl-top-sel}
\small
\resizebox{\textwidth}{!}{%
\begin{tabular}{llccccc}
\toprule
\textbf{\textit{Dataset}} & \textbf{Modelo} & \textbf{Taxa de aprend.} & \textbf{\textit{Weight Decay}} & \textbf{\textit{Dropout}} & \textbf{Dim. oculta} & \textbf{Lote} \\
\midrule
UCSB & LSTM & $3\cdot 10^{-3}$ & $10^{-4}$ & 0{,}3 & 256 & 4 \\
UCSB & BiLSTM & $3\cdot 10^{-3}$ & $10^{-4}$ & 0{,}3 & 512 & 4 \\
UCSB & GRU & $10^{-3}$ & $10^{-4}$ & 0{,}1 & 256 & 16 \\
UCSB & BiGRU & $3\cdot 10^{-3}$ & $10^{-4}$ & 0{,}3 & 256 & 8 \\
UCSB & SAC-LSTM & $10^{-2}$ & $10^{-4}$ & 0{,}3 & 64 & 4 \\
UCSB & MLP & $10^{-2}$ & $10^{-4}$ & 0{,}5 & 256 & 4 \\
\midrule
CR & LSTM & $3\cdot 10^{-2}$ & $10^{-4}$ & 0{,}3 & 128 & 32 \\
CR & BiLSTM & $10^{-2}$ & $10^{-4}$ & 0{,}3 & 64 & 32 \\
CR & GRU & $10^{-2}$ & $10^{-4}$ & 0{,}1 & 128 & 16 \\
CR & BiGRU & $3\cdot 10^{-3}$ & $10^{-4}$ & 0{,}5 & 256 & 64 \\
CR & SAC-LSTM & $3\cdot 10^{-3}$ & $10^{-4}$ & 0{,}5 & 64 & 64 \\
CR & MLP & $3\cdot 10^{-3}$ & $10^{-4}$ & 0{,}3 & 512 & 64 \\
\midrule
LA & LSTM & $3\cdot 10^{-4}$ & $10^{-4}$ & 0{,}1 & 128 & 16 \\
LA & BiLSTM & $3\cdot 10^{-4}$ & $10^{-4}$ & 0{,}1 & 512 & 16 \\
LA & GRU & $3\cdot 10^{-4}$ & $10^{-4}$ & 0{,}1 & 64 & 16 \\
LA & BiGRU & $3\cdot 10^{-4}$ & $10^{-4}$ & 0{,}1 & 64 & 16 \\
LA & SAC-LSTM & $3\cdot 10^{-4}$ & $10^{-4}$ & 0{,}1 & 512 & 32 \\
LA & MLP & $10^{-3}$ & $10^{-4}$ & 0{,}1 & 64 & 16 \\
\midrule
LG & LSTM & $3\cdot 10^{-3}$ & $10^{-4}$ & 0{,}5 & 256 & 16 \\
LG & BiLSTM & $3\cdot 10^{-3}$ & $10^{-4}$ & 0{,}1 & 256 & 16 \\
LG & GRU & $3\cdot 10^{-3}$ & $10^{-4}$ & 0{,}1 & 256 & 32 \\
LG & BiGRU & $3\cdot 10^{-3}$ & $10^{-4}$ & 0{,}5 & 512 & 32 \\
LG & SAC-LSTM & $10^{-3}$ & $10^{-4}$ & 0{,}1 & 512 & 32 \\
LG & MLP & $3\cdot 10^{-3}$ & $10^{-4}$ & 0{,}5 & 256 & 32 \\
\midrule
NHL & LSTM & $3\cdot 10^{-3}$ & $10^{-4}$ & 0{,}1 & 256 & 64 \\
NHL & BiLSTM & $10^{-3}$ & $10^{-4}$ & 0{,}5 & 512 & 16 \\
NHL & GRU & $10^{-3}$ & $10^{-4}$ & 0{,}3 & 256 & 16 \\
NHL & BiGRU & $3\cdot 10^{-3}$ & $10^{-4}$ & 0{,}3 & 256 & 32 \\
NHL & SAC-LSTM & $10^{-3}$ & $10^{-4}$ & 0{,}3 & 512 & 64 \\
NHL & MLP & $3\cdot 10^{-4}$ & $10^{-4}$ & 0{,}1 & 512 & 16 \\
\bottomrule
\end{tabular}%
}
\end{table}

A análise dos desempenhos aponta pontuações elevadas em todos os seis codificadores avaliados na validação cruzada. A BiGRU alcançou a maior pontuação média global nas partições de validação ($0{,}976$), com o melhor resultado registrado no conjunto UCSB ($0{,}933 \pm 0{,}061$), onde o volume amostral é mais restrito. A GRU e o MLP alcançaram as maiores pontuações no conjunto CR ($0{,}993 \pm 0{,}017$), enquanto o MLP obteve o maior índice no conjunto NHL ($0{,}973 \pm 0{,}020$). Todos os modelos alcançaram separação perfeita ($1{,}000 \pm 0{,}000$) no conjunto LA e índices superiores a $0{,}995$ no conjunto LG, evidenciando que as variações multiescala das medidas descritoras capturam de forma discriminativa a progressão tecidual associada à idade e ao tipo de tecido.

A distribuição dos hiperparâmetros selecionados evidencia a relevância da expansão da dimensão oculta. As dimensões ocultas ampliadas de $h = 256$ ou $h = 512$ foram selecionadas em 21 das 30 configurações ótimas (70\%), confirmando que espaços latentes mais amplos preservam com maior fidelidade as dinâmicas de conexão entre os raios. As taxas de aprendizado iguais ou superiores a $\eta = 3\cdot 10^{-3}$ prevaleceram em 18 configurações, atingindo até $3\cdot 10^{-2}$ no conjunto CR, sendo mantidas estáveis pela limitação da norma do gradiente. O tamanho de lote reduzido $B = 4$ foi o mais selecionado no conjunto UCSB em decorrência do volume restrito de imagens, enquanto tamanhos de $B = 16$ e $B = 32$ prevaleceram nos conjuntos com maior quantidade de amostras.

\section{Ablação da Via Visual}
\label{sec:ablacao-visual}

A via visual processa diretamente as imagens histopatológicas no formato de $256 \times 256$ pixels, obtidas pela preservação da proporção original com recorte central. Para aumentar a diversidade amostral durante o treinamento, foi aplicado aumento de dados rotacional composto por inversões horizontais e rotações ortogonais de 90 graus. A ablação da via visual foi conduzida com o propósito de avaliar o comportamento de diferentes abordagens de visão computacional e determinar as combinações ótimas de hiperparâmetros de cada modelo para a subsequente integração na arquitetura de dupla via. Foram avaliadas seis arquiteturas distribuídas em três famílias (\ref{fig:ablacao-vias}b):

\begin{itemize}
    \item a ResNet-50 \cite{he2016deep}, rede convolucional profunda com conexões residuais, empregada como referência padrão na classificação de tecidos histopatológicos \cite{litjens2017survey,histImgColorectal},
    \item a EfficientNet-B0 \cite{tan2019efficientnet}, arquitetura convolucional compacta baseada em escalonamento composto uniforme e blocos invertidos com mecanismo de atenção por canal \cite{kushwah2026evcnet,yousaf2024twostream},
    \item o ViT-Tiny \cite{dosovitskiy2020image}, modelo que processa a imagem como uma sequência unidimensional de projeções lineares de \textit{patches} por meio de camadas de autoatenção \cite{shamshad2023transformers,clinicalReview},
    \item o DeiT-Tiny \cite{touvron2021training}, transformador visual projetado com técnicas de destilação e regularização voltadas ao treinamento com menor dependência de grandes volumes de dados \cite{oliveira2025hybrid},
    \item a MobileViT-S \cite{mehta2022mobilevit}, arquitetura híbrida leve que intercala convoluções locais com blocos transformadores para equilibrar detalhamento celular e contexto amplo com baixo consumo paramétrico \cite{kabir2025semobilevit,ding2023himvit},
    \item a CoAtNet-0 \cite{dai2021coatnet}, modelo híbrido hierárquico que unifica convoluções nos primeiros estágios com autoatenção relativa nos estágios finais \cite{oliveira2025hybrid,zaidkilani2025dualstream}.
\end{itemize}

Todos os seis modelos foram inicializados com pesos pré-treinados na base ImageNet-1k e ajustados de ponta a ponta. A camada original de classificação foi substituída por uma cabeça linear adaptada à dimensão de saída de cada arquitetura, composta por camada linear, normalização por camadas, camada oculta de 128 neurônios com ativação ReLU, camada de \textit{dropout} e camada final para o número de classes de cada conjunto de dados.

O protocolo experimental seguiu a mesma divisão e metodologia adotadas na via topológica (\ref{fig:ablacao-treino}). A partição de teste final de 20\% foi mantida totalmente reservada e isolada do procedimento de busca. Sobre a partição de desenvolvimento de 80\%, foi executada uma validação cruzada estratificada em cinco partições ($k = 5$), com semente 42. Para cada modelo e para cada combinação de hiperparâmetros, foram executadas cinco rodadas nas quais quatro partições foram destinadas ao ajuste dos parâmetros e uma foi utilizada para a validação. A grade de hiperparâmetros explorada e as dimensões amostrais de cada conjunto de dados estão sintetizadas na Tabela~\ref{tab:espaco-busca-vis}.

\begin{table}[htb]
\centering
\caption{Dimensões amostrais e espaço de busca de hiperparâmetros estratificado por conjunto de dados na via visual.}
\label{tab:espaco-busca-vis}
\small
\resizebox{\textwidth}{!}{%
\begin{tabular}{lccccccc}
\toprule
\textbf{\textit{Dataset}} & \textbf{Desenv. (80\%)} & \textbf{Treino / Val. (rodada)} & \textbf{Lote ($B$)} & \textbf{Taxa do \textit{backbone} ($\eta_{bb}$)} & \textbf{Taxa da cabeça ($\eta_{hd}$)} & \textbf{Decaimento ($\lambda$)} & \textbf{\textit{Dropout} ($p$)} \\
\midrule
UCSB & 46  & 36 a 37 / 9 a 10   & $\{4, 8, 16\}$   & $\{10^{-5}, 3\cdot 10^{-5}, 10^{-4}, 3\cdot 10^{-4}\}$ & $\{10^{-4}, 3\cdot 10^{-4}, 10^{-3}\}$ & $10^{-4}$ & $\{0{,}1, 0{,}3, 0{,}5\}$ \\
CR   & 132 & 105 a 106 / 26 a 27 & $\{16, 32, 64\}$ & $\{10^{-5}, 3\cdot 10^{-5}, 10^{-4}, 3\cdot 10^{-4}\}$ & $\{10^{-4}, 3\cdot 10^{-4}, 10^{-3}\}$ & $10^{-4}$ & $\{0{,}1, 0{,}3, 0{,}5\}$ \\
LG   & 212 & 169 a 170 / 42 a 43 & $\{16, 32, 64\}$ & $\{10^{-5}, 3\cdot 10^{-5}, 10^{-4}, 3\cdot 10^{-4}\}$ & $\{10^{-4}, 3\cdot 10^{-4}, 10^{-3}\}$ & $10^{-4}$ & $\{0{,}1, 0{,}3, 0{,}5\}$ \\
NHL  & 299 & 239 a 240 / 59 a 60 & $\{16, 32, 64\}$ & $\{10^{-5}, 3\cdot 10^{-5}, 10^{-4}, 3\cdot 10^{-4}\}$ & $\{10^{-4}, 3\cdot 10^{-4}, 10^{-3}\}$ & $10^{-4}$ & $\{0{,}1, 0{,}3, 0{,}5\}$ \\
LA   & 422 & 337 a 338 / 84 a 85 & $\{16, 32, 64\}$ & $\{10^{-5}, 3\cdot 10^{-5}, 10^{-4}, 3\cdot 10^{-4}\}$ & $\{10^{-4}, 3\cdot 10^{-4}, 10^{-3}\}$ & $10^{-4}$ & $\{0{,}1, 0{,}3, 0{,}5\}$ \\
\bottomrule
\end{tabular}%
}
\end{table}

O treinamento foi conduzido com a função de perda de entropia cruzada e o otimizador AdamW. O decaimento de pesos foi fixado em $10^{-4}$ e a norma máxima dos gradientes foi limitada a 1,0. O processo foi limitado a 50 épocas, com parada antecipada estabelecida em 10 épocas consecutivas sem decréscimo da perda na partição de validação.

O desempenho de cada configuração foi mensurado pelo F1 macro médio nas cinco partições de validação interna. A combinação que atingiu a maior pontuação média foi definida como a configuração ótima de cada modelo em cada conjunto de dados. Os valores de F1 macro obtidos pelas melhores configurações de cada arquitetura estão reunidos na Tabela~\ref{tab:abl-vis-f1}.

\begin{table}[htb]
\centering
\caption{F1 macro (média $\pm$ desvio padrão entre as 5 partições de validação) da via visual isolada, na melhor combinação de hiperparâmetros de cada modelo. Em negrito, o maior valor médio de cada conjunto de dados.}
\label{tab:abl-vis-f1}
\small
\resizebox{\textwidth}{!}{%
\begin{tabular}{llccccc}
\toprule
\textbf{Família} & \textbf{Modelo} & \textbf{UCSB} & \textbf{CR} & \textbf{LA} & \textbf{LG} & \textbf{NHL} \\
\midrule
Convolucional & ResNet-50 & \pend & \pend & \pend & \pend & \pend \\
Convolucional & EfficientNet-B0 & \pend & \pend & \pend & \pend & \pend \\
\textit{Transformer} & ViT-Tiny & \pend & \pend & \pend & \pend & \pend \\
\textit{Transformer} & DeiT-Tiny & \pend & \pend & \pend & \pend & \pend \\
Híbrida & MobileViT-S & \pend & \pend & \pend & \pend & \pend \\
Híbrida & CoAtNet-0 & \pend & \pend & \pend & \pend & \pend \\
\bottomrule
\end{tabular}%
}
\end{table}

As combinações ótimas de hiperparâmetros responsáveis pelos maiores índices de F1 macro em cada conjunto de dados estão detalhadas na Tabela~\ref{tab:abl-vis-params}.

\begin{table}[htb]
\centering
\caption{Melhores combinações de hiperparâmetros obtidas na validação cruzada para cada modelo visual nos cinco conjuntos de dados.}
\label{tab:abl-vis-params}
\small
\resizebox{\textwidth}{!}{%
\begin{tabular}{lllccccc}
\toprule
\textbf{\textit{Dataset}} & \textbf{Família} & \textbf{Modelo} & $\eta_{bb}$ & $\eta_{hd}$ & \textbf{\textit{Dropout}} & \textbf{Lote} & \textbf{F1 macro} \\
\midrule
UCSB & Convolucional & ResNet-50 & \pend & \pend & \pend & \pend & \pend \\
UCSB & Convolucional & EfficientNet-B0 & \pend & \pend & \pend & \pend & \pend \\
UCSB & \textit{Transformer} & ViT-Tiny & \pend & \pend & \pend & \pend & \pend \\
UCSB & \textit{Transformer} & DeiT-Tiny & \pend & \pend & \pend & \pend & \pend \\
UCSB & Híbrida & MobileViT-S & \pend & \pend & \pend & \pend & \pend \\
UCSB & Híbrida & CoAtNet-0 & \pend & \pend & \pend & \pend & \pend \\
\midrule
CR & Convolucional & ResNet-50 & \pend & \pend & \pend & \pend & \pend \\
CR & Convolucional & EfficientNet-B0 & \pend & \pend & \pend & \pend & \pend \\
CR & \textit{Transformer} & ViT-Tiny & \pend & \pend & \pend & \pend & \pend \\
CR & \textit{Transformer} & DeiT-Tiny & \pend & \pend & \pend & \pend & \pend \\
CR & Híbrida & MobileViT-S & \pend & \pend & \pend & \pend & \pend \\
CR & Híbrida & CoAtNet-0 & \pend & \pend & \pend & \pend & \pend \\
\midrule
LA & Convolucional & ResNet-50 & \pend & \pend & \pend & \pend & \pend \\
LA & Convolucional & EfficientNet-B0 & \pend & \pend & \pend & \pend & \pend \\
LA & \textit{Transformer} & ViT-Tiny & \pend & \pend & \pend & \pend & \pend \\
LA & \textit{Transformer} & DeiT-Tiny & \pend & \pend & \pend & \pend & \pend \\
LA & Híbrida & MobileViT-S & \pend & \pend & \pend & \pend & \pend \\
LA & Híbrida & CoAtNet-0 & \pend & \pend & \pend & \pend & \pend \\
\midrule
LG & Convolucional & ResNet-50 & \pend & \pend & \pend & \pend & \pend \\
LG & Convolucional & EfficientNet-B0 & \pend & \pend & \pend & \pend & \pend \\
LG & \textit{Transformer} & ViT-Tiny & \pend & \pend & \pend & \pend & \pend \\
LG & \textit{Transformer} & DeiT-Tiny & \pend & \pend & \pend & \pend & \pend \\
LG & Híbrida & MobileViT-S & \pend & \pend & \pend & \pend & \pend \\
LG & Híbrida & CoAtNet-0 & \pend & \pend & \pend & \pend & \pend \\
\midrule
NHL & Convolucional & ResNet-50 & \pend & \pend & \pend & \pend & \pend \\
NHL & Convolucional & EfficientNet-B0 & \pend & \pend & \pend & \pend & \pend \\
NHL & \textit{Transformer} & ViT-Tiny & \pend & \pend & \pend & \pend & \pend \\
NHL & \textit{Transformer} & DeiT-Tiny & \pend & \pend & \pend & \pend & \pend \\
NHL & Híbrida & MobileViT-S & \pend & \pend & \pend & \pend & \pend \\
NHL & Híbrida & CoAtNet-0 & \pend & \pend & \pend & \pend & \pend \\
\bottomrule
\end{tabular}%
}
\end{table}

As configurações ótimas identificadas nas ablações da via topológica e da via visual fornecem os parâmetros fixos para a composição da arquitetura unificada de dupla via. O treinamento da rede de dupla via com os modelos selecionados e a avaliação do ganho preditivo frente às vias isoladas e aos métodos de referência no conjunto de teste independente são apresentados no Capítulo~\ref{chap:resultados}.
